% What is corvus
Corvus is a scientific workflow framework built with python with 
a focus on calculations of spectroscopy. It is written to facilitate
advanced calculations for novice users. 

Corvus is "property driven" which means that the user
specifies a property to calculate, and Corvus figures out 
how it can obtain that property, and what inputs are needed
from the user. This is done by building a dependency tree, with the root
given by the list of requested properties. The dependency tree is built
by checking the dependencies of (other properties needed to calculate) the 
requested property, then asking if these can be satisfied by any
of the available codes. This procedure is repeated until either all dependency
requirements are met, in which case a workflow is built and run. If all requirements
are not met, the user is notivied of any necessary user input.

Most of the workflows available to standard users are meant to facilitate calculations
using the FEFF10 code. These workflows include loops over specified parameters, configurational
averaging, RIXS, and creation of input files from multiple available structural input such as
.cif, .xyz, materials project id, or vasp output. 
